% Options for packages loaded elsewhere
\PassOptionsToPackage{unicode}{hyperref}
\PassOptionsToPackage{hyphens}{url}
%
\documentclass[
]{article}
\author{}
\date{}

\usepackage{amsmath,amssymb}
\usepackage{lmodern}
\usepackage{iftex}
\ifPDFTeX
  \usepackage[T1]{fontenc}
  \usepackage[utf8]{inputenc}
  \usepackage{textcomp} % provide euro and other symbols
\else % if luatex or xetex
  \usepackage{unicode-math}
  \defaultfontfeatures{Scale=MatchLowercase}
  \defaultfontfeatures[\rmfamily]{Ligatures=TeX,Scale=1}
\fi
% Use upquote if available, for straight quotes in verbatim environments
\IfFileExists{upquote.sty}{\usepackage{upquote}}{}
\IfFileExists{microtype.sty}{% use microtype if available
  \usepackage[]{microtype}
  \UseMicrotypeSet[protrusion]{basicmath} % disable protrusion for tt fonts
}{}
\makeatletter
\@ifundefined{KOMAClassName}{% if non-KOMA class
  \IfFileExists{parskip.sty}{%
    \usepackage{parskip}
  }{% else
    \setlength{\parindent}{0pt}
    \setlength{\parskip}{6pt plus 2pt minus 1pt}}
}{% if KOMA class
  \KOMAoptions{parskip=half}}
\makeatother
\usepackage{xcolor}
\IfFileExists{xurl.sty}{\usepackage{xurl}}{} % add URL line breaks if available
\IfFileExists{bookmark.sty}{\usepackage{bookmark}}{\usepackage{hyperref}}
\hypersetup{
  hidelinks,
  pdfcreator={LaTeX via pandoc}}
\urlstyle{same} % disable monospaced font for URLs
\usepackage{longtable,booktabs,array}
\usepackage{calc} % for calculating minipage widths
% Correct order of tables after \paragraph or \subparagraph
\usepackage{etoolbox}
\makeatletter
\patchcmd\longtable{\par}{\if@noskipsec\mbox{}\fi\par}{}{}
\makeatother
% Allow footnotes in longtable head/foot
\IfFileExists{footnotehyper.sty}{\usepackage{footnotehyper}}{\usepackage{footnote}}
\makesavenoteenv{longtable}
\usepackage{graphicx}
\makeatletter
\def\maxwidth{\ifdim\Gin@nat@width>\linewidth\linewidth\else\Gin@nat@width\fi}
\def\maxheight{\ifdim\Gin@nat@height>\textheight\textheight\else\Gin@nat@height\fi}
\makeatother
% Scale images if necessary, so that they will not overflow the page
% margins by default, and it is still possible to overwrite the defaults
% using explicit options in \includegraphics[width, height, ...]{}
\setkeys{Gin}{width=\maxwidth,height=\maxheight,keepaspectratio}
% Set default figure placement to htbp
\makeatletter
\def\fps@figure{htbp}
\makeatother
\setlength{\emergencystretch}{3em} % prevent overfull lines
\providecommand{\tightlist}{%
  \setlength{\itemsep}{0pt}\setlength{\parskip}{0pt}}
\setcounter{secnumdepth}{-\maxdimen} % remove section numbering
\ifLuaTeX
  \usepackage{selnolig}  % disable illegal ligatures
\fi

\begin{document}

\hypertarget{multicore-hardware-counter-correlation-study}{%
\section{Multicore Hardware Counter Correlation
Study}\label{multicore-hardware-counter-correlation-study}}

\hypertarget{executive-summary}{%
\subsection{Executive Summary}\label{executive-summary}}

This project collects Linux performance-counter measurements, aligns
them with workload metadata and, when available, synthetic phase labels,
and studies correlation structure to identify redundant indicators. The
main recommendation workflow prefers the normalized Spearman view
because raw counters often correlate simply because total activity rose
across a sample, not because two indicators capture the same underlying
behavior. On this machine, the generated artifact is based on PARSEC
workloads collected through \texttt{parsecmgmt}.

\hypertarget{objective}{%
\subsection{Objective}\label{objective}}

Determine which hardware indicators are highly correlated or redundant
so a later phase-classification and phase-prediction pipeline can rely
on a smaller, cheaper, and more portable counter set.

\hypertarget{experimental-setup}{%
\subsection{Experimental Setup}\label{experimental-setup}}

\begin{itemize}
\tightlist
\item
  CPU vendor/model: \texttt{GenuineIntel} /
  \texttt{Intel(R)\ Xeon(R)\ CPU\ E5-2680\ v2\ @\ 2.80GHz}
\item
  Architecture: \texttt{x86\_64}
\item
  Logical CPUs: \texttt{20}
\item
  Sockets: \texttt{2}
\item
  Cores per socket: \texttt{10}
\item
  Perf version: \texttt{perf\ version\ 5.14.0-611.47.1.el9\_7.x86\_64}
\item
  Intel PCM available: \texttt{False}
\item
  CPU family/model id: \texttt{6} / \texttt{62}
\item
  Strict core PMU study ready: \texttt{True}
\item
  Uncore IMC study ready: \texttt{False}
\end{itemize}

\hypertarget{workload-generation}{%
\subsection{Workload Generation}\label{workload-generation}}

\begin{itemize}
\tightlist
\item
  Suites in this generated artifact: \texttt{parsec}.
\item
  PARSEC available on this host: \texttt{True}. SPEC CPU2017 available
  on this host: \texttt{False}.
\item
  Workloads executed:
  \texttt{blackscholes,\ bodytrack,\ canneal,\ fluidanimate,\ freqmine}.
\item
  Thread counts executed: \texttt{1,\ 2,\ 4,\ 8}.
\item
  Repetitions per workload/thread combination: \texttt{3}.
\item
  Target collection interval: \texttt{10\ ms}.
\item
  PARSEC input sizes used: \texttt{test}.
\item
  Observed perf-wrapped interval collection wall time:
  \texttt{1594.8\ ms\ median\ (1193.0\ min,\ 2550.9\ max)}.
\item
  Affinity policy: Pinned with \texttt{taskset} to the first N online
  CPUs for the requested thread count.
\item
  \texttt{parsec} workload generation: standard PARSEC applications and
  kernels launched through \texttt{parsecmgmt} with taskset-based CPU
  affinity and the configured PARSEC input set.
\end{itemize}

\hypertarget{event-mapping-table}{%
\subsection{Event Mapping Table}\label{event-mapping-table}}

\begin{longtable}[]{@{}
  >{\raggedright\arraybackslash}p{(\columnwidth - 6\tabcolsep) * \real{0.25}}
  >{\raggedright\arraybackslash}p{(\columnwidth - 6\tabcolsep) * \real{0.25}}
  >{\raggedright\arraybackslash}p{(\columnwidth - 6\tabcolsep) * \real{0.25}}
  >{\raggedright\arraybackslash}p{(\columnwidth - 6\tabcolsep) * \real{0.25}}@{}}
\toprule
\begin{minipage}[b]{\linewidth}\raggedright
Family
\end{minipage} & \begin{minipage}[b]{\linewidth}\raggedright
Selected Event
\end{minipage} & \begin{minipage}[b]{\linewidth}\raggedright
Supported
\end{minipage} & \begin{minipage}[b]{\linewidth}\raggedright
Note
\end{minipage} \\
\midrule
\endhead
Branch instructions & br\_inst\_retired.all\_branches & Yes & Using
preferred host event \\
Branch mispredictions & br\_misp\_retired.all\_branches & Yes & Using
preferred host event \\
Cycles & cpu\_clk\_unhalted.thread & Yes & Using preferred host event \\
Instructions retired & inst\_retired.any & Yes & Using preferred host
event \\
L1 data loads & mem\_uops\_retired.all\_loads & Yes & Using preferred
host event \\
L1 data stores & mem\_uops\_retired.all\_stores & Yes & Using preferred
host event \\
LLC misses & longest\_lat\_cache.miss & Yes & Using preferred host
event \\
LLC references & longest\_lat\_cache.reference & Yes & Using preferred
host event \\
Memory read bandwidth & uncore\_imc/cas\_count\_read/ & Yes & Preferred
Intel uncore IMC event (system-wide) \\
Memory write bandwidth & uncore\_imc/cas\_count\_write/ & Yes &
Preferred Intel uncore IMC event (system-wide) \\
Off-core / DRAM demand reads & offcore\_requests.demand\_data\_rd & Yes
& Using preferred host event \\
Resource stalls & resource\_stalls.any & Yes & Using preferred host
event \\
Total memory bandwidth & derived\_total\_memory\_bandwidth & Yes &
Derived from system-wide Intel uncore IMC read+write \\
\bottomrule
\end{longtable}

\begin{itemize}
\tightlist
\item
  Confident counters retained for analysis:
  \texttt{Branch\ instructions,\ Branch\ mispredictions,\ Cycles,\ Instructions\ retired,\ L1\ data\ loads,\ L1\ data\ stores,\ LLC\ misses,\ LLC\ references,\ Off-core\ /\ DRAM\ demand\ reads,\ Resource\ stalls}.
\item
  Confident counters exposed on this host:
  \texttt{Branch\ instructions,\ Branch\ mispredictions,\ Cycles,\ Instructions\ retired,\ L1\ data\ loads,\ L1\ data\ stores,\ LLC\ misses,\ LLC\ references,\ Memory\ read\ bandwidth,\ Memory\ write\ bandwidth,\ Off-core\ /\ DRAM\ demand\ reads,\ Resource\ stalls,\ Total\ memory\ bandwidth}.
\item
  Generic fallback counters excluded from analysis: \texttt{none}.
\item
  Unsupported counter families excluded from analysis:
  \texttt{Floating-point\ arithmetic,\ L2\ misses}.
\item
  Missing required strict-study families: \texttt{none}.
\item
  Uncore readiness note:
  \texttt{System-wide\ perf\ is\ blocked\ on\ this\ host.\ perf\_event\_paranoid=2;\ lower\ it\ or\ grant\ perf\ capability\ for\ uncore\ collection.}.
\end{itemize}

\hypertarget{dataset-summary}{%
\subsection{Dataset Summary}\label{dataset-summary}}

\begin{longtable}[]{@{}ll@{}}
\toprule
Statistic & Value \\
\midrule
\endhead
Interval rows & 9409 \\
Aggregate rows & 0 \\
Runs merged & 60 \\
Manifest runs & 60 \\
Stale raw runs excluded & 213 \\
Observed raw counter columns & 10 \\
Confident counter columns kept & 10 \\
Generic fallback columns excluded & 0 \\
Columns removed as constant & 0 \\
Strict core PMU study ready & Yes \\
Uncore IMC study ready & No \\
Uncore blocked by host policy & Yes \\
Observed system-wide uncore columns & 0 \\
\bottomrule
\end{longtable}

\hypertarget{preprocessing-and-derived-metrics}{%
\subsection{Preprocessing and Derived
Metrics}\label{preprocessing-and-derived-metrics}}

The preprocessing stage handles missing values with median imputation
after logging missingness, removes constant columns, creates raw and
winsorized views, and derives CPI, IPC, MPKI-style cache/branch/off-core
metrics, and FP intensity metrics. Raw-counter and normalized-counter
views are intentionally separated because raw counts are strongly
scale-sensitive. For normalized and derived analysis views, samples with
non-positive interval duration, instructions retired, or cycles are
excluded to avoid startup/shutdown artifacts.

\begin{itemize}
\tightlist
\item
  Analysis-ready rows kept: \texttt{9306} of \texttt{9409}.
\item
  Dropped rows: \texttt{103}.
\item
  Drop reasons:
  \texttt{non\ positive\ cycles:\ 100,\ non\ positive\ instructions\ retired:\ 100,\ non\ positive\ interval\ duration\ ms:\ 0}.
\end{itemize}

\hypertarget{global-correlation-results}{%
\subsection{Global Correlation
Results}\label{global-correlation-results}}

Raw counters show which indicators move together at the absolute-count
level, but these relationships may largely reflect workload intensity.
Normalized counters reduce that effect and are the better basis for
redundancy decisions. This Xeon host exposes Intel uncore IMC PMUs, but
system-wide collection is currently blocked by host policy, so the
present artifact reflects confident core PMU analysis only.

\hypertarget{raw-pearson-heatmap}{%
\subsubsection{Raw Pearson Heatmap}\label{raw-pearson-heatmap}}

\begin{figure}
\centering
\includegraphics{../tables/raw/pearson_heatmap.png}
\caption{Raw Pearson Heatmap}
\end{figure}

\hypertarget{raw-spearman-heatmap}{%
\subsubsection{Raw Spearman Heatmap}\label{raw-spearman-heatmap}}

\begin{figure}
\centering
\includegraphics{../tables/raw/spearman_heatmap.png}
\caption{Raw Spearman Heatmap}
\end{figure}

\hypertarget{normalized-pearson-heatmap}{%
\subsubsection{Normalized Pearson
Heatmap}\label{normalized-pearson-heatmap}}

\begin{figure}
\centering
\includegraphics{../tables/normalized/pearson_heatmap.png}
\caption{Normalized Pearson Heatmap}
\end{figure}

\hypertarget{normalized-spearman-heatmap}{%
\subsubsection{Normalized Spearman
Heatmap}\label{normalized-spearman-heatmap}}

\begin{figure}
\centering
\includegraphics{../tables/normalized/spearman_heatmap.png}
\caption{Normalized Spearman Heatmap}
\end{figure}

\hypertarget{clustered-heatmap}{%
\subsubsection{Clustered Heatmap}\label{clustered-heatmap}}

\begin{figure}
\centering
\includegraphics{../tables/normalized/clustered_heatmap.png}
\caption{Clustered Heatmap}
\end{figure}

\hypertarget{dendrogram}{%
\subsubsection{Dendrogram}\label{dendrogram}}

\begin{figure}
\centering
\includegraphics{../tables/normalized/dendrogram.png}
\caption{Dendrogram}
\end{figure}

\hypertarget{stratified-results}{%
\subsection{Stratified Results}\label{stratified-results}}

Stratified analysis is emitted by workload and thread count when enough
samples exist. Without explicit phase labels, workload-local stability
is the main check on whether redundancy patterns persist across
applications.

\hypertarget{top-positive-and-negative-correlations}{%
\subsubsection{Top Positive and Negative
Correlations}\label{top-positive-and-negative-correlations}}

\begin{figure}
\centering
\includegraphics{../tables/normalized/top_correlations.png}
\caption{Top Positive and Negative Correlations}
\end{figure}

\hypertarget{multicollinearity-results}{%
\subsection{Multicollinearity Results}\label{multicollinearity-results}}

Variance inflation factors complement pairwise correlation: a counter
can have moderate pairwise correlations but still be redundant when
explained by several others together. VIF is therefore best used
alongside correlation-threshold grouping instead of as a replacement.

\hypertarget{vif-bar-chart}{%
\subsubsection{VIF Bar Chart}\label{vif-bar-chart}}

\begin{figure}
\centering
\includegraphics{../tables/normalized/vif.png}
\caption{VIF Bar Chart}
\end{figure}

\hypertarget{feature-selection-recommendations}{%
\subsection{Feature Selection
Recommendations}\label{feature-selection-recommendations}}

\hypertarget{g01}{%
\subsubsection{G01}\label{g01}}

\begin{itemize}
\tightlist
\item
  Keep: \texttt{Cycles\ /\ ms}
\item
  Drop: Instructions retired / ms
\item
  Members: Cycles / ms, Instructions retired / ms
\item
  Reason: Selected for higher portability/interpretability and lower
  missingness; max within-group \textbar rho\textbar=0.976
\end{itemize}

\hypertarget{g02}{%
\subsubsection{G02}\label{g02}}

\begin{itemize}
\tightlist
\item
  Keep: \texttt{LLC\ references\ /\ KI}
\item
  Drop: Off-core / DRAM demand reads / KI
\item
  Members: LLC references / KI, Off-core / DRAM demand reads / KI
\item
  Reason: Selected for higher portability/interpretability and lower
  missingness; max within-group \textbar rho\textbar=0.976
\end{itemize}

The full machine-readable recommendation table remains available in
\texttt{../tables/recommendations.csv}.

\hypertarget{limitations}{%
\subsection{Limitations}\label{limitations}}

\begin{itemize}
\tightlist
\item
  Raw counters can correlate simply because all counts grow with
  activity, sample duration, or thread count.
\item
  Vendor-specific PMU names differ substantially; unsupported counters
  are logged and left blank.
\item
  Intel PCM support is best-effort because deployments vary by binary
  name, permissions, and CLI syntax.
\item
  Per-core task-local attribution is limited without privileged
  system-wide PMU access.
\item
  Intel uncore IMC bandwidth on this host requires system-wide perf
  access; when \texttt{perf\_event\_paranoid} blocks that path, the
  report explicitly omits those system-wide metrics.
\end{itemize}

\hypertarget{next-steps}{%
\subsection{Next Steps}\label{next-steps}}

\begin{itemize}
\tightlist
\item
  Validate the recommended reduced set against a downstream phase
  classifier or predictor.
\item
  Repeat the study on additional Intel, AMD, and Arm systems to compare
  portability.
\item
  Extend the stratified analysis with application-specific phases from
  PARSEC or SPEC CPU2017 when those suites are available.
\item
  Lower \texttt{perf\_event\_paranoid} or grant equivalent perf
  capability on this Xeon host if you want IMC bandwidth collected
  alongside the strict core PMU set.
\end{itemize}

\hypertarget{appendix}{%
\subsection{Appendix}\label{appendix}}

\begin{itemize}
\tightlist
\item
  Raw indicators analyzed:
  \texttt{Instructions\ retired,\ Cycles,\ Branch\ instructions,\ Branch\ mispredictions,\ L1\ data\ loads,\ L1\ data\ stores,\ LLC\ references,\ LLC\ misses,\ Off-core\ /\ DRAM\ demand\ reads,\ Resource\ stalls}
\item
  Derived indicators analyzed: \texttt{6}
\item
  Counter families defined in the project:
  \texttt{instructions\_retired,\ cycles,\ branch\_instructions,\ branch\_mispredictions,\ l1d\_loads,\ l1d\_stores,\ l2\_misses,\ llc\_references,\ llc\_misses,\ offcore\_demand\_data\_reads,\ fp\_arithmetic,\ resource\_stalls,\ memory\_read\_bandwidth,\ memory\_write\_bandwidth,\ total\_memory\_bandwidth}
\end{itemize}

\hypertarget{how-to-interpret-the-results}{%
\subsection{How to interpret the
results}\label{how-to-interpret-the-results}}

\begin{itemize}
\tightlist
\item
  Two counters are likely redundant when they remain highly correlated
  in the normalized Spearman view and also appear in the same high-VIF
  neighborhood.
\item
  Raw and normalized views can disagree because raw counts mix true
  behavior with overall workload activity and duration scaling.
\item
  Use correlation to find pairwise or group redundancy, then use VIF to
  see whether an indicator adds information once the rest of the set is
  considered.
\item
  Choose a minimal portable counter set by preferring common counters
  with low missingness and clear meaning, then dropping vendor-specific
  counters only when a more portable representative preserves the same
  signal.
\end{itemize}

\end{document}
